% !PRINTED 658
\weberpage{658}

Dann würde sich aus (20) eine Kongruenz ergeben:

\[ y_0 + y_1\tau + y_2\tau^2 + \cdots + y_{n-1}\tau^{n-1} \equiv 0 \quad (\mathrm{mod}\ z-c), \]

was nach dem Satz 5. nicht möglich ist.

Nach § 174, (17) und § 176, (9) ist

\weq{27}{\begin{aligned}
N F'(\tau) &= \pm \Delta(1, \tau, \tau^2, \ldots, \tau^{n-1}) \\
&= \pm U^2 \Delta(\omega_1, \omega_2, \ldots, \omega_n) \\
&= \pm U^2 \Delta;
\end{aligned}}

also ist

\weq{28}{N_a F'(\tau) = \Delta}

\gesperrt{die Körperdiskriminante, die hiernach durch (24) in lineare Faktoren
zerlegt ist.}

Die gewonnenen Resultate benutzen wir noch zum Beweis des folgenden
Satzes:

7. \gesperrt{Ist eine ganze Funktion in $\Omega$}

\weq{29}{\omega = x_1\omega_1 + x_2\omega_2 + \cdots + x_n\omega_n}

\gesperrt{durch jedes in $z-c$ aufgehende Primfunktional teilbar, so
ist $S(\omega)$ durch $z-c$ teilbar.}

Die $x_1, x_2, \ldots, x_n$ sind hier ganze Funktionen in $Z$. Die
Funktion $\omega$ geht aus $\tau$ hervor durch die Substitution

\weq{30}{t_1 = x_1,\ t_2 = x_2,\ \ldots,\ t_n = x_n.}

Ist also $\omega$ durch $\pi$ teilbar, und $\omega_1 \equiv c_1,
\ldots, \omega_n \equiv c_n \pmod{\pi}$, so ist

\[ x_1c_1 + x_2c_2 + \cdots + x_nc_n \equiv 0 \quad (\mathrm{mod}\ \pi) \]

und folglich auch, als Funktion in $Z$, durch $z-c$ teilbar. Sind
also $\tau_1, \tau_2, \ldots, \tau_m$ die konstanten Funktionale,
denen $\tau$ nach den Primfaktoren $\pi_1, \pi_2, \ldots, \pi_m$ von
$(z-c)$ kongruent ist, und gehen diese durch die Substitution (30) in
$\tau_1^0, \tau_2^0, \ldots, \tau_m^0$ über, so sind alle diese
Funktionen durch $(z-c)$ teilbar [weil sie n. V. durch einen
Primteiler von $(z-c)$ teilbar sind]. Nun ist nach (22)

\weq{31}{-S(\tau) \equiv e_1\tau_1 + e_2\tau_2 + \cdots + e_m\tau_m,}

und mithin

\[ -S(\omega) \equiv e_1\tau_1^0 + e_2\tau_2^0 + \cdots + e_m\tau_m^0 \equiv 0 \quad (\mathrm{mod}\ z-c), \]

wie zu beweisen war.

\begin{center}
§ 185. \textbf{\gesperrt{Die gebrochenen Funktionen in $\Omega$ und die
Taylorsche Entwickelung.}}
\end{center}

Eine gebrochene Funktion $\eta$ im Körper $\Omega$ kann auf
unendlich viele Arten durch Multiplikation mit einer ganzen Funktion
$\nu$ in eine ganze Funktion $\mu$ verwandelt werden. Es ist dann
%CONT>>
